% ==================================================
% RELATED WORK
% ==================================================

\section{Related Work}

\subsection{Sequential Recommendation Systems}

Sequential recommendation has become increasingly important as researchers recognize that user behavior is inherently temporal and sequential in nature.

\subsubsection{RNN-Based Approaches}

Early works used RNNs to capture sequential patterns:

\begin{itemize}
    \item Hidasi et al. introduced GRU4REC, using gated recurrent units for session-based recommendations
    \item Extensions like improved GRU4REC added data augmentation and techniques to address data sparsity
\end{itemize}

\subsubsection{Attention-Based Methods}

The introduction of self-attention in sequential recommendation improved performance and interpretability:

\begin{itemize}
    \item SASRec (Self-Attentive Sequential Recommendation) proposed by Kang et al. uses multi-head self-attention
    \item Subsequent works extended this with graph structures and improved architectures
\end{itemize}

\subsubsection{Graph-Enhanced Approaches}

Recent methods combine graphs with sequential models:

\begin{itemize}
    \item gSASRec augments SASRec with item relationships through graph embeddings
    \item Other works explore heterogeneous graphs for multi-type interactions
\end{itemize}

\subsection{Sequential Recommendation in E-Learning}

E-learning systems have unique characteristics that affect recommendation design:

\begin{itemize}
    \item Learning outcomes depend on prerequisites and learning sequences
    \item Pedagogical constraints must be considered (prerequisites, difficulty progression)
    \item Students have diverse learning goals and backgrounds
    \item Course structures impose constraints on valid recommendation sequences
\end{itemize}

Few works specifically target sequential recommendation for e-learning platforms, making this an important research direction.

\subsection{Dataset Studies}

Research on the MARS dataset and similar e-learning datasets:

\begin{itemize}
    \item MARS provides opportunities to study real student behavior patterns
    \item Prior analysis has identified key characteristics like resource popularity skew
    \item Temporal patterns in student access reveal learning rhythm variations
\end{itemize}

\subsection{Performance Metrics}

Common metrics for evaluating sequential recommendation systems include:

\begin{enumerate}
    \item \textbf{Recall@K:} Proportion of relevant items in top-K recommendations
    \item \textbf{NDCG@K:} Normalized discounted cumulative gain considering ranking quality
    \item \textbf{MRR:} Mean reciprocal rank of the first relevant item
    \item \textbf{Hit Rate:} Frequency with which correct item appears in top-K
\end{enumerate}

\subsection{Gap Analysis}

Despite the growing attention to sequential recommendation and e-learning systems, there remains a gap in:

\begin{itemize}
    \item Specialized models for educational domain constraints
    \item Comprehensive evaluation on real e-learning datasets like MARS
    \item Practical deployment considerations for production systems
\end{itemize}

This thesis addresses these gaps by proposing and evaluating a gSASRec-based system specifically for the MARS e-learning context.

\newpage
